\chapter{Requirements}\label{ch:requirements}
% Source: docs/requirements.md (v0.2). Keep IDs identical so tests and report agree.

The requirements were written in the first two weeks of the project from the approved proposal
and the two scope decisions (ADR 0001 and ADR 0002), and were reviewed with the advisor on
23 October 2026. Each requirement has an identifier that is used unchanged in the code, the tests
and this report. The goals H1--H5 are the objectives of Chapter~\ref{ch:intro}.

\section{Stakeholders}
Table~\ref{tab:stakeholders} lists who uses the system or depends on its results.

\begin{table}[ht]
  \centering
  \small
  \caption{Stakeholders and their needs.}
  \label{tab:stakeholders}
  \begin{tabularx}{\textwidth}{@{}lX@{}}
    \toprule
    Role & Needs \\
    \midrule
    Applicant (employed or student) & know which documents are needed, upload them, see problems before applying \\
    Administrator & maintain requirements and expert notes, review detected source changes before they go live \\
    Researcher (project author) & generate labelled synthetic files, run and compare approaches, measure results \\
    Advisor and course coordinator & follow progress, agree on success criteria, assess the outcome \\
    \bottomrule
  \end{tabularx}
\end{table}

\section{Functional requirements}
Table~\ref{tab:fr} lists the functional requirements. \emph{Must} is required for the semester's
outcome, \emph{Should} is planned, \emph{Could} is done only if time allows. The last column gives
the state on 20 November 2026; Chapter~\ref{ch:implementation} describes what was built.

{\small
\begin{longtable}{@{}p{1.1cm}>{\raggedright\arraybackslash}p{6.0cm}p{1.0cm}p{1.3cm}p{1.5cm}p{1.2cm}@{}}
  \caption{Functional requirements.}\label{tab:fr}\\
  \toprule
  ID & Requirement & Goal & Priority & Semester & State \\
  \midrule
  \endfirsthead
  \toprule
  ID & Requirement & Goal & Priority & Semester & State \\
  \midrule
  \endhead
  \bottomrule
  \endfoot
  \reqid{FR-01} & The applicant selects destination, visa type, profile and travel dates. & H1 & Must & Fall & done \\
  \reqid{FR-02} & A personalised checklist; each item states how it is submitted and links to its official source. & H1, H4 & Must & Fall & done \\
  \reqid{FR-03} & A discovery page shows Schengen membership, appointment and biometrics information and official links. & H1, H4 & Should & Fall & done \\
  \reqid{FR-04} & Users register and log in; roles are applicant and administrator. & -- & Must & Fall & done \\
  \reqid{FR-05} & The applicant creates, lists and deletes application files. & -- & Must & Fall & done \\
  \reqid{FR-06} & The applicant tracks the application steps. & H1 & Should & Fall & done \\
  \reqid{FR-07} & Upload per checklist item (PDF, JPEG, PNG), including from a phone camera. & H2 & Must & Fall & done \\
  \reqid{FR-08} & The applicant enters or confirms the key fields of each document. & H2 & Must & Fall & done \\
  \reqid{FR-09} & Rule checks detect missing, outdated and inconsistent documents. & H2 & Must & Fall & done \\
  \reqid{FR-10} & Every finding shows the rule, the documents and fields involved, the official source and a severity. & H2, H4 & Must & Fall & done \\
  \reqid{FR-11} & Deleting a document or an application file removes the stored files. & -- & Must & Fall & done \\
  \reqid{FR-12} & The system classifies the type of an uploaded document. & H2 & Must & Spring & -- \\
  \reqid{FR-13} & Fields are extracted with OCR and a vision-language model, with a confidence per field; low-confidence fields are confirmed by the applicant. & H2 & Must & Fall, Spring & ongoing \\
  \reqid{FR-14} & Cross-document checks cover all pilot document pairs. & H2 & Must & Spring & partly \\
  \reqid{FR-15} & An assistant answers from official sources and expert notes, cites its passages and gives no definitive answer without a source. & H4 & Must & Spring & -- \\
  \reqid{FR-16} & Official pages are captured periodically; changes are detected, summarised, linked to requirements and approved by an administrator. & H3 & Must & Spring & partly \\
  \reqid{FR-17} & After an approved change, affected applicants are notified. & H3 & Should & Spring & -- \\
  \reqid{FR-18} & Administrators manage requirements and expert notes in the web application. & H3, H4 & Should & Spring & -- \\
  \reqid{FR-19} & Long-running checks report their status to the browser. & -- & Should & Spring & -- \\
  \reqid{FR-20} & The web application can be installed on a phone. & -- & Could & Spring & -- \\
  \reqid{FR-21} & Applicants can message an expert. & -- & Could & Spring & -- \\
  \reqid{FR-22} & Research tooling generates labelled synthetic files, runs an LLM-only baseline and computes the metrics. & H5 & Must & Spring & partly \\
\end{longtable}
}

The checklist (\reqid{FR-02}) states whether an original, a copy or both are submitted, and the
translation requirement of Norway; a notarisation label is not modelled because neither pilot
source asks for one. \reqid{FR-14} is partly done because the name and passport-number
comparisons already run across all pilot documents on the entered fields. \reqid{FR-16} is partly
done because the pages are captured weekly and compared by hash; detecting meaningful changes and
the approval step belong to spring. \reqid{FR-13} moved partly into the fall semester after the
advisor meeting of 20 November 2026: the prototype reads the fields of the synthetic documents
with a hosted model, and the applicant confirms them; OCR, confidence-based routing and a model
that can run on the project's own server follow in spring. \reqid{FR-22} is partly done because
two labelled demo files and the command that scores the rule engine on them exist
(Chapter~\ref{ch:evaluation}).

\section{Non-functional requirements}
Table~\ref{tab:nfr} lists the non-functional requirements. Two of them shape the design more than
the others. Privacy (\reqid{NFR-01}): the application files hold copies of passports, bank
statements and payslips, which are personal data under the Turkish data-protection law (KVKK);
Chapter~\ref{ch:ethics} describes the measures. Languages (\reqid{NFR-10}): the applicants read
Turkish, the official sources are in English, so the interface is in both languages and quoted
passages stay in their original language.

\begin{table}[ht]
  \centering
  \small
  \caption{Non-functional requirements.}
  \label{tab:nfr}
  \begin{tabularx}{\textwidth}{@{}p{3.4cm}X@{}}
    \toprule
    ID & Requirement \\
    \midrule
    \reqid{NFR-01} Privacy & consent before the first upload; documents encrypted at rest; limited retention; deletion on request; fields not needed for a check masked before any external AI call; no real personal documents in development, tests or experiments \\
    \reqid{NFR-02} Security & salted password hashes; token-based authentication; access to one's own files only; secrets only in \texttt{.env}; the database not exposed to the network \\
    \reqid{NFR-03} Transparency & every requirement, finding and answer links to its source; every result page states that the system prepares and does not predict a decision \\
    \reqid{NFR-04} Usability & usable on a 375\,px wide phone screen and on desktop; status and findings understandable without technical knowledge \\
    \reqid{NFR-05} Performance & rule checks within 2 seconds for one application file on the development machine; extraction time per document measured in spring \\
    \reqid{NFR-06} Cost & model usage within USD 20 per month; the cost of every evaluation run logged \\
    \reqid{NFR-07} Reproducibility & the system starts with Docker Compose following the README; pinned dependencies; lint and tests in CI on every pull request \\
    \reqid{NFR-08} Maintainability & every rule has unit tests; decisions recorded as ADRs; code and repository documents in English \\
    \reqid{NFR-09} Robustness & uploads validated for type and size (at most 10\,MB); liveness and readiness checks \\
    \reqid{NFR-10} Languages & Turkish and English interface, selectable at any time; all interface text from translation files; quoted passages in their original language \\
    \bottomrule
  \end{tabularx}
\end{table}

\section{Constraints}
One developer at about 15 hours per week over two semesters; an API budget of at most USD 20 per
month; evaluation on synthetic documents only, so results may not fully transfer to real
applications. Official pages can change without notice; the pilot relies on the sources described
in Chapter~\ref{ch:background}. The pilot covers one main destination, a second national layer,
one visa category and two profiles.

\section{Success criteria}\label{sec:success}
Agreed with the advisor on 23 October 2026: every result is a comparison against a reference or
an alternative, never an isolated score. Results are reported as measured; no score is promised
in advance.

\paragraph{End of fall (prototype).}
\begin{itemize}
  \item An applicant of each pilot profile completes the flow from profile to checklist, upload,
    field entry and findings.
  \item The generated checklist is compared item by item with a checklist compiled by hand from
    the official sources; every item and every rule cites its source.
  \item On the demo application files, the findings of the rule engine are compared with the
    injected error labels: every injected error of the supported types is reported, and the
    error-free file gets no finding.
  \item The checklist of a pilot profile is compared with the answer of a general-purpose chat
    assistant to the same question: missing items, wrong items and missing citations.
  \item The language model that reads the fields is compared with the same model used alone on
    the same demo files: labelled errors found, false alarms, time and cost (added after the
    advisor meeting of 20 November 2026).
  \item The system starts by following the README; all tests pass in CI.
\end{itemize}

\paragraph{End of spring (complete project).}
\begin{itemize}
  \item A labelled synthetic data set (about 200 files, more than 1,000 documents) with a frozen
    test set.
  \item Hybrid and LLM-only approaches compared per error type with precision, recall, F1,
    false-alarm rate, time and cost per file (AS1).
  \item Human reviewers compared with the hybrid system on a subset of the test set (AS2).
  \item Change detection compared with labelled page versions against a plain text-diff baseline
    (AS3).
  \item Assistant answers compared with LLM-only answers: citation coverage and unsupported
    statements.
\end{itemize}
